\section{Summation Notation}

\begin{definition}[Summation]
Let \(f\) be a function defined on the integers.
For integers \(a\) and \(b\) with \(a \le b\), the expression

\[
\sum_{k=a}^{b} f(k)
\]

denotes the sum

\[
f(a) + f(a+1) + f(a+2) + \cdots + f(b).
\]

The symbol \(\sum\) (the Greek capital letter sigma) represents the
operation of summation.
\end{definition}

\subsection*{Elements of the Notation}

The summation expression

\[
\sum_{k=a}^{b} f(k)
\]

contains several components:

\begin{itemize}
\item \(\sum\) — the summation symbol.
\item \(k\) — the \emph{index of summation}.
\item \(a\) — the \emph{lower limit} (starting index).
\item \(b\) — the \emph{upper limit} (ending index).
\item \(f(k)\) — the \emph{summand}, the expression being added.
\end{itemize}

Thus the index \(k\) runs through the integer values

\[
a, a+1, a+2, \dots, b.
\]

\begin{remark}
The index of summation is a \emph{dummy variable}.  
Its name is irrelevant, and the expression

\[
\sum_{k=a}^{b} f(k)
\]

is identical to

\[
\sum_{i=a}^{b} f(i).
\]
\end{remark}

\subsection*{Basic Properties of Summation}

For functions \(f\) and \(g\) and constants \(a,b,c\),

\begin{align*}
\sum_{k=m}^{n} (f(k) + g(k))
&=
\sum_{k=m}^{n} f(k)
+
\sum_{k=m}^{n} g(k)
\\[6pt]
\sum_{k=m}^{n} c\,f(k)
&=
c \sum_{k=m}^{n} f(k)
\end{align*}

These properties express the linearity of summation.

\subsection*{Splitting Sums}

A sum over an interval may be divided into smaller sums:

\[
\sum_{k=a}^{c} f(k)
=
\sum_{k=a}^{b} f(k)
+
\sum_{k=b+1}^{c} f(k)
\]

whenever \(a \le b < c\).

\subsection*{Common Examples}

\[
\sum_{k=1}^{n} 1 = n
\]

\[
\sum_{k=1}^{n} k = \frac{n(n+1)}{2}
\]

\[
\sum_{k=1}^{n} k^2 = \frac{n(n+1)(2n+1)}{6}
\]

These formulas frequently appear in discrete mathematics and analysis.

\subsection*{Nested Summations}

Summations may be nested in a way analogous to nested loops.

For example,

\[
\sum_{i=1}^{n} \sum_{j=1}^{m} a_{ij}
\]

means

\[
\sum_{i=1}^{n}
\big(
a_{i1} + a_{i2} + \cdots + a_{im}
\big).
\]

The inner sum is evaluated first.

\begin{remark}
Nested sums behave similarly to iterated integrals.
Under appropriate conditions the order of summation may be exchanged:

\[
\sum_{i=1}^{n}\sum_{j=1}^{m} a_{ij}
=
\sum_{j=1}^{m}\sum_{i=1}^{n} a_{ij}.
\]
\end{remark}

\subsection*{Changing the Index}

It is often useful to change the summation index.

Example:

\[
\sum_{k=0}^{n} f(k)
=
\sum_{j=1}^{n+1} f(j-1)
\]

Such substitutions are frequently used to simplify expressions
or align sums with known formulas.

\subsection*{Empty Sums}

If the upper bound is smaller than the lower bound, the sum contains no terms.
By convention,

\[
\sum_{k=a}^{b} f(k) = 0
\qquad
\text{when } a > b.
\]

This convention simplifies many algebraic identities.